\chapter*{摘要}
\addcontentsline{toc}{chapter}{摘要}

本文档为 \textbf{ebbackup 备份内核}（目录 \filepath{engine\_cpp/}）与 \textbf{运维质量闭环}（bench L1--L7、compact/GC、232 项 gtest、C API ABI v12）的\textbf{完整扩写版}工程技术手册。目标篇幅 \textbf{120+ 页}（含附录与索引），较 v1.0 摘要版显著加深：每一核心子系统均给出字节级布局、算法步骤、字段表、版本对照与实测基准。

\textbf{Part I} 自底向上解析备份内核：Content ID 恒为 \texttt{SHA256(content)}；增量路径 HCRBO + CFI 锚点 reuse；提交点耐久性要求 manifest 写入 superblock 前必须 \texttt{ChunkStore::Flush()}（fsync）。Pipeline 自 \wave{4}（\versiontag{0.8.0}）支持全局 chunk 队列；\wave{5}（\versiontag{0.9.0}）修复 ChunkStore 并发并引入 Stage 3.1/3.2 分层门禁；\versiontag{0.9.4} Sprint 4 在流式 CDC 上增加 Phase B seg1 bulk 与 Phase A \texttt{FastCdcSlice} opt-in 路由。

\textbf{Part II} 解析 bench 硬门禁、C API ABI、compact/GC 运维与 CI 矩阵，定义产品可交付质量标准，且\textbf{不改变} Part I 中的 Content ID 与 commit-point 语义。

\vspace{1em}
\noindent\textbf{读者对象：}
\begin{itemize}[nosep]
  \item 备份引擎/存储系统开发者；
  \item 数据保护与 dedup 研究人员；
  \item 对接 \filepath{eb\_backup.h} C API 或 \texttt{eb} CLI 的集成工程师；
  \item 负责 CI/bench 门禁与性能回归的 SRE/QA。
\end{itemize}

\vspace{1em}
\noindent\textbf{文档约定：}
\begin{itemize}[nosep]
  \item 源码路径均相对于仓库根 \filepath{E:/recoveryProjects/}；
  \item 流程图使用 TikZ；代码清单摘录自头文件/实现，行号以文档编写时 commit \texttt{4a7ac94}（\versiontag{0.9.4}）为准；
  \item 吞吐单位默认 \textbf{SI MB/s}（$10^6$ B/s），与 \filepath{docs/reference/PERF\_BASELINE.md} 一致；
  \item 版本演进详见 \filepath{CHANGELOG.md}、\filepath{docs/VERSION\_HISTORY.md}。
\end{itemize}

\vspace{1em}
\noindent\textbf{文档结构概览：}
\begin{longtable}{@{}lp{10cm}@{}}
\toprule
Part & 章节 \\
\midrule
Part I 内核 & 概述、架构、不变量、BackupEngine、FastCDC/HCRBO、Digest/Crypto、ChunkStore/EbPack、Manifest/Snapshot、Pipeline v4、流式 CDC、恢复/校验、审计/RAR、CLI/daemon \\
Part II 运维 & bench L1--L7、C API ABI、compact/GC、测试/CI \\
附录 & 源码索引、术语表、环境变量、manifest v4/HXID 字节参考、Wave 演进、统计字段索引 \\
\bottomrule
\end{longtable}

\chapter*{修订历史}
\addcontentsline{toc}{chapter}{修订历史}
\begin{longtable}{@{}llp{9cm}@{}}
\toprule
日期 & 版本 & 说明 \\
\midrule
2026-07-07 & v1.0 & 首版：Part I 内核 12 章 + Part II 运维 4 章 + 附录（摘要级） \\
2026-07-07 & v1.1 & \textbf{完整扩写}：01/02/03/17--20/99 深度重写；目标 120+ 页；对齐 \versiontag{0.9.4} ABI v12、232 gtest、Stage 3.1 floors \\
\bottomrule
\end{longtable}
